% answers_short.tex —— 2011 年·上海春季卷 · 简短答案速查
% 由 tools/make_short_answers.py 自 answers.tex 的 \gksolhead/\gkanswer 提取，
% 并与 2011/2011_数学真题_上海春季_答案要点.md 逐题比对后生成，勿手改（改源 MD 后重跑）。

\begin{gknotice}
\textbf{使用说明}：本节只给最终答案便于快速核对，完整推导见解析版。答案为编写者独立推导，\textbf{非官方标准答案}。
\end{gknotice}

\gkpart{填空题}
\begin{itemize}[leftmargin=2.6em]
  \item[\textbf{1.}] $(2,+\infty)$。
  \item[\textbf{2.}] $\{x\mid1\leqslant x\leqslant2\}$。
  \item[\textbf{3.}] $\dfrac{\sqrt{22}}{11}$。
  \item[\textbf{4.}] $1$。
  \item[\textbf{5.}] $\arcsin\dfrac13$。
  \item[\textbf{6.}] $20$。
  \item[\textbf{7.}] $\dfrac\pi{12}$。
  \item[\textbf{8.}] $-7$。
  \item[\textbf{9.}] $\dfrac{x^2}9+\dfrac{y^2}4=1$。
  \item[\textbf{10.}] $2$。
  \item[\textbf{11.}] $7$。
  \item[\textbf{12.}] $168$。
  \item[\textbf{13.}] $\dfrac\pi3$。
  \item[\textbf{14.}] $\dfrac{-1-\sqrt5\pm\sqrt{10-2\sqrt5}\mathrm i}4$ 或 $\dfrac{-1+\sqrt5\pm\sqrt{10+2\sqrt5}\mathrm i}4$（写出其中一个即可）。
\end{itemize}

\gkpart{选择题}
\begin{center}\begin{tabular}{ccccc}
\toprule
题号 & 15 & 16 & 17 & 18 \\
\midrule
答案 & C & A & D & B \\
\bottomrule
\end{tabular}\end{center}

\gkpart{解答题}
\begin{itemize}[leftmargin=2.6em]
  \item[\textbf{19.}] 最小正周期 $\pi$；$x\in\left[0,\dfrac\pi2\right]$ 时的最大值为 $\sqrt2$（$x=\dfrac\pi8$ 时取得）。
  \item[\textbf{20.}] 表面积约 $87.96\ \mathrm{cm}^2$，体积约 $57.80\ \mathrm{cm}^3$。
  \item[\textbf{21.}] (1) $0$；(2) $k_{P_1P_2}-k_{P_2P_3}+k_{P_3P_4}-\cdots+(-1)^{n-1}k_{P_nP_1}=\dfrac{1+(-1)^{n-1}}2$（$P_1=P(2,1)$），即 $n$ 为偶数时为 $0$，$n$ 为奇数时为 $1$
  \item[\textbf{22.}] (1) 证明见解析；(2) $f(x)=-x^2$（说明见解析）；(3) $f(x)=\begin{cases}x^2,&x\geqslant1,\\ x,&x<1.\end{cases}$
  \item[\textbf{23.}] (1) $x_1=4.74$，$x_2=4.67$，$x_3=4.65$；$x_{n+1}<x_n$；(2) 证明见解析；(3) 由（2）的误差估计取 $n=16$；若用计算器对给定 $x_0$ 逐项验算，$n=10$ 即可保证（参考答案同时给出两种估计）
\end{itemize}
